\section{ベイズ推定と判断}
\label{s27_bayesian3}

\subsection{授業内容}
\subsubsection{科目の中でのこのコマの位置づけ}
\begin{tabular}[t]{cccccc}
	記述統計[3] & 確率[13] & モデル[15] & 推測・推定[8] & 検定[6] & 実践[5]
\end{tabular}

\subsubsection{概要}
ベイズ推定による結果の解釈は，分布の区間として直接的な解釈ができるため，帰無仮説検定で用いていたような「差があると言って間違っている確率」という持って回った言い回しを必要としない。
ここではまず区間推定を行い，判断基準としての帰無仮説・対立仮説の導入という従来法による群間差の判断原理を再確認する。つづいてベイズ法による区間推定と，区間を使ってどのような判断が可能なのか(HDI区間，ROPEによる判断)を考える。
また，帰無仮説検定が二つのモデルの比較であったことを思い出し，ベイズ流のモデル比較であるベイズファクターについて考える。
ベイズファクターの計算には周辺尤度が必要であるが，サヴェージ・ディッキー法をつかうとその簡便的な計算が可能であることにも触れる。

\subsubsection{コマ主題細目(箇条書き)}
\begin{description}
	\item[ベイズ法と帰無仮説検定] ベイズ法による事後分布を用いた区間推定は，信用区間ではなく確信区間であり，区間が分布の情報を含んだものであること，母数に対する確信度であることを確認する。一方で，モーメント法での推定は母数を含んでいると判断すると，当たっている確率としての幅であり，判断そのものは当たっているか外れているかの0/1判断である。同様のことが帰無仮説検定における判断にも言える。改めて帰無仮説検定の判断ロジックを確認しておく。

	      \begin{flushright}
		      $\to$ \citet{Toyoda201606}のP.42--43.
	      \end{flushright}

	\item[区間による判定] ベイズ法を用いて群間差を評価する方法として，区間推定の結果を用いるROPEを使う方法について解説する。区間が母数の分布を含んでいるために判断がしやすい反面，明確な基準がなく，またROPEをどのように設定するのかについてはドメイン知識が必要になることが難点であるが，これはそもそも従来法でも考えるべき問題であった。またベイズ法による事後分布は一つに限定されるので，多重比較の問題が生じないことにも触れる。
	      \begin{flushright}
		      $\to$ \citet{Maeda2017}のP.303-334
	      \end{flushright}

	\item[ベイズファクター] ベイズ法によるもう一つの判断方法は，帰無仮説検定と同じくモデル比較の観点から考えることである。ベイズファクターは周辺尤度の比からいずれのモデルがデータに支持されているかを表す指標であり，帰無仮説モデルと対立仮説モデルをこれで比較することによって，どちらがよりデータに適合したモデルであるかを評価することができる。特徴として，帰無仮説のようなNull仮説を積極的に支持することができる点がある。そもそも帰無仮説モデルが対立仮説モデルに比べて非常に限定的であること，ネストされたモデルであることを指摘し，このような場合はサヴェージ・ディッキー法で比較できることを解説する。
	      \begin{flushright}
		      $\to$ \citet{Lee201709}のP.88--103.
	      \end{flushright}

\end{description}
\subsubsection{キーワード}
\begin{itemize}
	\item 標本分布と事後分布
	\item 区間推定と点推定
	\item ベイズファクター
	\item サヴェージ・ディッキー法
\end{itemize}
\subsection{授業情報}
\paragraph{コマの展開方法}
講義
\paragraph{標準シラバスにおける位置づけ}
科目番号5；心理学統計法；(1)心理学で用いられる統計手法；J より高度な記述や推定を目指して
\subsubsection{予習・復習課題}
\paragraph{予習}帰無仮説検定の前段階である，母数の区間推定について改めて理解しておく。また帰無仮説検定が推定結果からある種の判断を行うためのツールであったこと，科学的探究のためには必ずしも0/1の判断結果に陥る必要がないことを確認しておく。
\paragraph{復習}帰無仮説検定の時と違って，どのような結論の出し方や解釈の仕方になるのか，しっかり理解しておこう。統計ツールとしてRのほかにJASPが導入される。JASPは2021年12月現在でもバージョン1.0に満たないソフトウェアではあるが，ベイズ法と従来法を並列的に活用できるツールであるので，興味があれば自分のPCにダウンロードして試してみるのも良いだろう。
